% TOP_PROBLEM: {"id": 5100028, "title": "Elliptic-billiard invariant k_{502}", "category_id": 11, "category": {"id": 11, "name": "dynamical_systems", "display_name": "Dynamical Systems", "description": "Problems about long-term behavior of deterministic systems, Hamiltonian dynamics, and periodic orbits.", "slug": "dynamical-systems", "order_index": 11, "created_at": "2026-07-31T15:26:25.671Z"}, "status": "partially_solved"}
% FIRST_POSTED: null
% ENTRY_KIND: statement_only
% Imported from: batch01.tex; UnsolvedMath ID 5100028
\documentclass[11pt]{article}
\usepackage[margin=1in]{geometry}
\usepackage{amsmath,amssymb,mathtools}
\usepackage{fontspec,unicode-math}
\setmainfont{Latin Modern Roman}
\setsansfont{Latin Modern Sans}
\setmathfont{Latin Modern Math}
\usepackage{xurl,hyperref}
\hypersetup{colorlinks=true,linkcolor=blue,urlcolor=blue}
\newcommand{\sref}[2]{\hyperlink{source:#1:#2}{[S#2]}}
\newcommand{\cataloguescope}[1]{\paragraph{Mathematical question.}#1}
\begin{document}
% UnsolvedMath ID 5100028; source catalogue code AMR-050-0028
\setcounter{section}{5100027}
\section{Elliptic-billiard invariant k\_\{502\}}\label{5100028}
{\small\sffamily UnsolvedMath ID 5100028 · Dynamical Systems}
\subsection{Definitions and mathematical statement}

\paragraph{Shared notation.}$\mathbb N=\{1,2,\ldots\}$, and $\mathbb Z,\mathbb Q,\mathbb R,\mathbb C$ denote the integers, rationals, reals and complex numbers. Any other conventions are specified locally.


Fix the ellipse $E:x^2/a^2+y^2/b^2=1$, where $a>b>0$, with center $O=(0,0)$ and foci $f_1,f_2=(\pm\sqrt{a^2-b^2},0)$. Fix a confocal elliptic caustic inside $E$ admitting a Poncelet family of $N$-periodic billiard polygons $P=(P_1,\ldots,P_N)$, with $N\geq3$. Their vertices lie on $E$, their sides are tangent to the caustic, and the incoming and outgoing sides obey the billiard reflection law. Indices are cyclic. The outer polygon $P'$ has vertices given by intersections of consecutive tangent lines to $E$ at the $P_i$; the inner polygon $P''$ has vertices at the contacts of the sides of $P$ with the caustic. All constancy assertions concern variation within this fixed family.
For any ordered polygon $R=(R_i)_{i=1}^N$, use the signed area
\[S(R)=\frac12\sum_{i=1}^N R_i\times R_{i+1},\qquad (x,y)\times(u,v)=xv-yu.\]
Write $A=S(P)$, $A'=S(P')$ and $A''=S(P'')$. A pedal polygon $\operatorname{ped}_M(R)$ has vertices the perpendicular feet from $M$ to the lines containing the sides of $R$. An antipedal polygon $\operatorname{ant}_M(R)$ has vertices at consecutive intersections of the lines through $R_i$ perpendicular to $R_i-M$. Define
\[A_M=S(\operatorname{ped}_M(P)),\quad A'_M=S(\operatorname{ped}_M(P')),\quad A_M^*=S(\operatorname{ant}_M(P)).\]
Primes or double primes on analogous quantities refer to $P'$ or $P''$. For $R$, its vertex centroid and signed area centroid are
\[C_0(R)=\frac1N\sum_i R_i,\qquad C_2(R)=\frac{1}{6S(R)}\sum_i(R_i\times R_{i+1})(R_i+R_{i+1}).\]
Constructions and quotients are understood wherever their defining intersections and denominators exist.
If $R$ has interior angles $\theta_i(R)$, its Steiner centroid of curvature in the source convention is
\[K(R)=\frac{\sum_i\sin(2\theta_i(R))R_i}{\sum_i\sin(2\theta_i(R))}.\]
Put $K=K(P)$, $K'=K(P')$, $K''=K(P'')$, and $A_K=S(\operatorname{ped}_K(P))$, $A'_{K'}=S(\operatorname{ped}_{K'}(P'))$, $A''_{K''}=S(\operatorname{ped}_{K''}(P''))$.
\paragraph{Statement recorded in the supplied source.}
Prove that for odd $N$, $A'/A'_{K'}$ is constant over the family (invariant $k_{502}$). \sref{5100028}{2}
\subsection{Short English statement}
For odd periods, prove that the signed area of the outer polygon divided by the signed area of its pedal with respect to its own Steiner curvature centroid stays constant.
\subsection{Sources}
\begin{description}
\item[\hypertarget{source:5100028:1}{S1}] ULAM.AI and the UnsolvedMath Contributors, UnsolvedMath: A Curated Collection of Open Mathematics Problems, record 5100028 (AMR-050-0028). CC BY 4.0. \url{https://huggingface.co/datasets/ulamai/UnsolvedMath}.
\item[\hypertarget{source:5100028:2}{S2}] Dan Reznik, Ronaldo Garcia and Jair Koiller, \emph{Eighty New Invariants in the Elliptic Billiard}, arXiv:2004.12497v11 (29 October 2020), \S2, equation (4), \S3.6 and Table 6, pp. 3, 6--8. \url{https://arxiv.org/abs/2004.12497v11}.
\end{description}
\label{end:5100028}
\end{document}
